\appendix

\section{Proof of Theorem~\ref{thm:level}}\label{app:proof}

We follow the proof of~\cite[Theorem~5.1]{BS23}. It suffices to consider deterministic provers. Let $\mathcal P^*$ be a deterministic prover that makes the verifier accept with probability $\varepsilon>\varepsilon_0$. In an accepting transcript we write $\vz=\vz^{(0)}+b\vz^{(1)}$, $\ve=\ve^{(0)}+b\ve^{(1)}$, $\tH_i=G_1\that_i$ and $h_i=G_w\hat{\mathbf h}_i$ for the values recomposed from the last message. Since $\norm{\vw'}\le\beta'$,
\begin{equation}\label{eq:app-norms}
\norm{(\vz,\ve)}\le\norm{(\vz^{(0)},\ve^{(0)})}+b\,\norm{(\vz^{(1)},\ve^{(1)})}\le(b+1)\beta',\qquad\norm{(\that,\vd_1)}\le\beta',\qquad\norm{(\hat{\mathbf h},\vd_2)}\le\beta' .
\end{equation}
Whenever we say that the prover is bound to a value, two accepting transcripts in which the value differs yield a solution to one of the instances of $\nfsis$ in the statement of the theorem; we argue as if this did not happen.

\paragraph{Step 1: outer commitments.}
The message $\vu_1$ is the same in all transcripts. By the third row of~\eqref{eq:next} and~\eqref{eq:app-norms}, $(\that,\vd_1)$ is an opening of norm $\beta'$ of the rounded commitment $\vu_1$ under the key $B$. By \cref{lem:binding}(i) the prover is bound to $\that$, hence to the $\tH_i$. In the same way the fourth row binds the prover to $\hat{\mathbf h}$, hence to the $h_i$, in all accepting transcripts that agree on the challenges sent before $\vu_2$, that is, on all challenges except $c_1,\dots,c_r$.

\paragraph{Step 2: weak openings.}
As in~\cite{BS23,NS22}, for each $i$ the extractor obtains two accepting transcripts that differ only in the challenge $c_i\ne c_i'$. It repeats the following until it succeeds: it runs $\mathcal P^*$ with fresh challenges until it finds an accepting transcript, and then runs it up to $N=\lceil1/p_0\rceil$ times with all challenges of this transcript except a fresh $c_i'$, where $p_0=\varepsilon/2-1/\lvert\Ch\rvert>0$; it succeeds if one of these runs is accepting with $c_i'\ne c_i$, and otherwise starts again with a new first transcript. Call the challenges other than $c_i$ of a transcript \emph{heavy} if $\mathcal P^*$ succeeds with probability at least $\varepsilon/2$ over the choice of $c_i$. The challenges of a random accepting transcript are heavy with probability at least $1/2$, because the transcripts that are accepting and not heavy have total probability less than $\varepsilon/2$. If they are heavy, each of the $N$ runs succeeds with probability at least $p_0$, and one of them does with probability at least $1-(1-p_0)^N\ge1-1/e$. An iteration thus succeeds with probability at least $(1-1/e)/2$ and costs $1/\varepsilon+N$ runs in expectation, so the expected number of runs for one $i$ is $O(1/p_0)=O\bigl((\varepsilon-2/\lvert\Ch\rvert)^{-1}\bigr)$. Let $(\vz,\ve)$ and $(\vz',\ve')$ be the two responses, $\bar c_i=c_i-c_i'$ and
\[
(\vs^*_i,\vl^*_i)=\bar c_i^{-1}\bigl((\vz,\ve)-(\vz',\ve')\bigr)\in\Rq^n\times\Rq^\kappa .
\]
Subtracting the first row of~\eqref{eq:next} for the two transcripts, in which the $\tH_j$ coincide by Step~1, gives
\begin{equation}\label{eq:app-weak}
A\vs^*_i-\vl^*_i=2^{\rin}\tH_i\qquad\text{and}\qquad\norm{\bar c_i(\vs^*_i,\vl^*_i)}\le2(b+1)\beta' .
\end{equation}
Thus $(\vs^*_i,\vl^*_i,\bar c_i)$ is a weak opening of $\tH_i$ of norm $2(b+1)\beta'$. It does not depend on the transcripts from which it was computed: by \cref{lem:binding}(iii), two different weak openings of this norm of the same $\tH_i$ yield a solution of $\nfsis_{\kappa,n,4T\cdot2(b+1)\beta'}$, which is the problem for $A$ in the statement of the theorem, and the $\tH_i$ are the same in all accepting transcripts by Step~1. In particular the vector $\vs^*=(\vs^*_1,\dots,\vs^*_r)$ is determined by the first message $\vu_1$, before any challenge is chosen, although the extractor can compute it only at the end. It is fixed from now on. We show that it is a witness for the norm bound $\sqrt{128/30}\,\beta$ by bounding the probability that a transcript is accepting if it is not.

\paragraph{Step 3: the form of the response.}
Fix $i_0$. For an accepting transcript let $\mathbf y=(\vz,\ve)-c_{i_0}(\vs^*_{i_0},\vl^*_{i_0})$. By the first row of~\eqref{eq:next} and~\eqref{eq:app-weak}, $[A\mid-I]\,\mathbf y=2^{\rin}\sum_{i\ne i_0}c_i\tH_i$, and by~\eqref{eq:app-norms} and~\eqref{eq:app-weak}
\[
\norm{\bar c_{i_0}\mathbf y}\le\opn{\bar c_{i_0}}\norm{(\vz,\ve)}+\opn{c_{i_0}}\norm{\bar c_{i_0}(\vs^*_{i_0},\vl^*_{i_0})}\le2T(b+1)\beta'+2T(b+1)\beta' .
\]
In two accepting transcripts that agree on all challenges except $c_{i_0}$, the vectors $\mathbf y$ are thus weak openings of norm $4T(b+1)\beta'$ of the same rounded commitment $\sum_{i\ne i_0}c_i\tH_i$ with the same $\bar c_{i_0}$. By \cref{lem:binding}(ii) and the hardness of $\nfsis_{\kappa,n,8T(b+1)\beta'}$ they coincide: in all accepting transcripts that differ only in $c_{i_0}$ we have $\vz=\mathbf y_z+c_{i_0}\vs^*_{i_0}$, where $\mathbf y_z$, the first component of $\mathbf y$, does not depend on $c_{i_0}$.

\paragraph{Step 4: partial evaluations.}
For challenges $\rho$ let $h^*_i=\ip{\vE}{\vs^*_i}$. We claim that an accepting transcript with $h_i\ne h^*_i$ for some $i$ occurs with probability at most $1/\lvert\Ch\rvert$. The $h_i$, the $h_i^*$ and $\vE$ are determined by the challenges other than $c_1,\dots,c_r$; let $i_0$ be the smallest index with $h_{i_0}\ne h^*_{i_0}$ and fix all challenges except $c_{i_0}$. By Step~3 the second row of~\eqref{eq:next} reads
\[
c_{i_0}\bigl(h^*_{i_0}-h_{i_0}\bigr)=\sum_{i\ne i_0}c_ih_i-\ip{\vE}{\mathbf y_z}
\]
in every accepting transcript, and the right-hand side does not depend on $c_{i_0}$. If the equation held for two values of $c_{i_0}$, their difference, which is invertible, would annihilate $h^*_{i_0}-h_{i_0}\ne0$. Hence at most one of the $\lvert\Ch\rvert$ values of $c_{i_0}$ leads to an accepting transcript.

\paragraph{Step 5: sum-check.}
Let $\vPhi$ and $b$ be as in~\eqref{eq:agg}. For $j\le\mu$ let $g_j^*$ be the polynomial of step~4 of \cref{sec:level} computed from $\vs^*$ and $\rho_1,\dots,\rho_{j-1}$, and let $g_j$ be the polynomial that the verifier reconstructs. Then $g_1^*(0)+g_1^*(1)=\Theta(\ip{\vPhi}{\vs^*})$, $g^*_{j+1}(0)+g^*_{j+1}(1)=g_j^*(\rho_j)$ and $g^*_\mu(\rho_\mu)=\Theta(\sum_i\Phi'_ih^*_i)$, while $g_1(0)+g_1(1)=b$, $g_{j+1}(0)+g_{j+1}(1)=g_j(\rho_j)$ and $g_\mu(\rho_\mu)=v$. Consider an accepting transcript with $h_i=h_i^*$ for all $i$ and $\Theta(\ip{\vPhi}{\vs^*})\ne b$. Then $g_1\ne g_1^*$, and $g_\mu(\rho_\mu)=g^*_\mu(\rho_\mu)$ by the fifth row of~\eqref{eq:next}. Let $j$ be the largest index with $g_j\ne g_j^*$. Then $g_j(\rho_j)=g_j^*(\rho_j)$, either because $j=\mu$ or because $g_{j+1}=g_{j+1}^*$. Both polynomials have degree at most $2$ and are determined before $\rho_j$ is chosen, so this happens with probability at most $2/\lvert\FF\rvert$ for each $j$, and such transcripts occur with probability at most $2\mu/\lvert\FF\rvert$.

\paragraph{Step 6: aggregation.}
Let $\hat y_k\in\FF$ be the lifted values used by the verifier, with $\tau_0(\hat y_k)=\ip{\omega^{(k)}}{\mathbf p}$. The vectors $\vs^*$ and $\vpsi_k$ and the values $\hat y_k$ are fixed before $\alpha$, $\xi$ and $\nu$ are chosen. By \cref{lem:agg}, applied to the constraints of the statement together with the $\FF$-constraints $\Theta(\ip{\vpsi_k}{\vs^*})=\hat y_k$, the probability that $\Theta(\ip{\vPhi}{\vs^*})=b$ although one of these constraints is violated is $1/\lvert\FF\rvert$.

\paragraph{Step 7: projection.}
Suppose that $\Theta(\ip{\vpsi_k}{\vs^*})=\hat y_k$ for all $k\le\kq$. Applying $\tau_0$ and using \cref{lem:theta}(iii) and~\eqref{eq:ct} gives $\sum_j\omega^{(k)}_j\bigl(\ip{\pi_j}{\vecc(\vs^*)}-p_j\bigr)=0$ in $\Z_q$ for all $k$. If $\mathbf p\ne\Pi\,\vecc(\vs^*)\bmod q$, this happens with probability $q^{-\kq}$ over the choice of the $\omega^{(k)}$, which are sent after $\mathbf p$.

\paragraph{Step 8: norm.}
If $\mathbf p=\Pi\,\vecc(\vs^*)\bmod q$ and $\norm{\vs^*}>\sqrt{128/30}\,\beta$, then by \cref{lem:jl} with $b=\sqrt{128/30}\,\beta\le q/125$ the verifier's check $\norm{\mathbf p}\le\sqrt{128}\,\beta=\sqrt{30}\,b$ passes with probability at most $\varepsilon_{\mathrm{JL}}$ over the choice of $\Pi$, which is sent after $\vu_1$.

\paragraph{Conclusion.}
The constraint vectors of the statement vanish on the padding, and removing coordinates does not increase the norm; so if the padded vector $\vs^*$ satisfies the constraints and the norm bound, its first $m$ coordinates are a witness, and it suffices to show the former. Suppose that $\vs^*$ violates a constraint of the statement or has norm larger than $\sqrt{128/30}\,\beta$. Then every accepting transcript belongs to one of the events of Steps~4 to~8: either some $h_i$ differs from $h_i^*$; or the aggregated claim is false for $\vs^*$; or it is true although a constraint of the statement or a lifted value is wrong; or all constraints hold and all lifted values are right, in which case the norm of $\vs^*$ is too large and $\mathbf p$ either is not its projection or is. The probabilities of these events add up to $1/\lvert\Ch\rvert+(2\mu+1)/\lvert\FF\rvert+q^{-\kq}+\varepsilon_{\mathrm{JL}}<\varepsilon_0<\varepsilon$, a contradiction. By Step~2 the extractor runs $\mathcal P^*$ an expected number of $O\bigl(r/(\varepsilon-2/\lvert\Ch\rvert)\bigr)\le O\bigl(r/(\varepsilon-\varepsilon_0)\bigr)$ times.

For the terminal level, replace Step~1 by the observation that the $\tH_i$ and $h_i$ are sent before the challenges on which they must not depend, and~\eqref{eq:app-norms} by $\norm{(\vz,\ve)}\le\bar\gamma$, which is~\eqref{eq:terminal}. \qed

\section{The Size Model}\label{app:model}

\paragraph{Ranks.}
For a norm $B$ the rank is the smallest $\kappa$ with $2\sqrt{\kappa d\log_2q\log_2\delta(\mathsf b)}>\log_2B$, see~\eqref{eq:mr}.

\paragraph{Coding.}
Vectors of small integers are coded with an arithmetic coder under a discretised Gaussian model.
\begin{lemma}\label{lem:coding}
Let $\sigma>0$. There is a prefix-free code that encodes every $\mathbf v\in\Z^N$ with $\norm{\mathbf v}^2\le N\sigma^2$ in at most $N\log_2\bigl((\sigma\sqrt{2\pi}+1)\sqrt e\bigr)+2$ bits.
\end{lemma}
\begin{proof}
Let $P(x)=e^{-x^2/2\sigma^2}/Z$ for $x\in\Z$. Arithmetic coding with the model $P$ uses at most $-\sum_i\log_2P(v_i)+2$ bits~\cite{WNC87}. Here $Z=\sum_xe^{-x^2/2\sigma^2}$ is at most $\sigma\sqrt{2\pi}+1$, and
\[
-\sum_i\log_2P(v_i)=N\log_2Z+\frac{\norm{\mathbf v}^2}{2\sigma^2\ln2}\le N\log_2Z+N\log_2\sqrt e .\qedhere
\]
\end{proof}
The bound depends only on the norm that the verifier checks anyway. We write $H(\sigma)=\log_2((\sigma\sqrt{2\pi}+1)\sqrt e)$; without coding the model uses $\log_2(12\sigma)$ bits per coefficient as in~\cite{BS23}. A high part $\hi_\varrho$ takes $\log_2(q/2^\varrho+2)$ bits per coefficient.

\paragraph{A recursive level.}
Let $\gamma^2=\tau\beta^2$ and $\gamma_e^2=\tau\,r\kappa d\,4^{\rin}/12$, the expected squared norms of $\vz$ and $\ve$ by \cref{lem:mean} for low parts that are uniform; for the bound of \cref{sec:completeness} both are doubled. With $\sigma_z=\gamma/\sqrt{nd}$ the base is $b=\lfloor(\sqrt{12}\,\sigma_z)^{1/2}\rceil$, the inner commitments have $t_1=\lfloor(\log_2q-\rin)/\log_2b\rceil$ digits in base $b_1=\lceil(q/2^{\rin})^{1/t_1}\rceil$, and the partial evaluations $t_w=\lfloor\log_2q/\log_2b\rceil$ digits in base $b_w=\lceil q^{1/t_w}\rceil$. Digits and low parts are counted with the variance of a uniform distribution, as in~\cite{BS23}:
\[
\beta'^2=\frac{2\gamma^2}{b^2}+\frac{b_1^2}{12}\,r\kappa t_1d+\frac{b_w^2}{12}\,rt_wd+\Bigl(\frac{b^2}{12}\,\kappa d+\frac{\gamma_e^2}{b^2}\Bigr)+\frac{4^{\rout}}{12}\,2\kappa_1d .
\]
The number $\rin$ of dropped bits is the largest for which $r\kappa d\,4^{\rin}/12$ is at most a fraction of $\beta^2$ that is part of the search, and $\rout$ the largest for which the last term is at most a fraction $\eta\in\{0.1,0.3\}$ of the others. Since $\kappa$ and $\kappa_1$ depend on $\beta'$, they are determined by fixed-point iteration. The level contributes
\[
2\kappa_1d\log_2\bigl(q/2^{\rout}+2\bigr)+256\,H\bigl(\beta/\sqrt2\bigr)+\kq(f-1)\log_2q+2\mu f\log_2q+512
\]
bits, and the next witness has $m'$ ring elements as in~\eqref{eq:mprime}. If $\vz$ is not decomposed or the partial evaluations are sent in the clear, the corresponding terms change in the obvious way.

\paragraph{The terminal level.}
It contributes $r\kappa d\log_2(q/2^{\rin}+2)+rd\log_2q+nd\,H(\sigma_z)$ bits for the inner commitments, the partial evaluations and $\vz$, plus the projection, the lifting, the sum-check and the seeds as above; $\kappa$ is chosen for $8T\bar\gamma$ in the proof-based and for $2\bar\gamma$ in the conventional setting, where $\bar\gamma^2=\gamma^2+\gamma_e^2$, and all $\rin\le16$ are tried.

\paragraph{Guaranteed parameters.}
In the guaranteed setting of \cref{sec:completeness} the variances above are replaced by squared maximal absolute values and the bounds on $\vz$ and $\ve$ by the joint bound~\eqref{eq:joint}; the norm bound of the next level is $B'$ from~\eqref{eq:bprime}. Since the verifier of the next level accepts projections for $\sqrt{1+t}\,B'$, the ranks are chosen for $2\varsigma\sqrt{1+t}\,B'$ and $8T(b+1)\varsigma\sqrt{1+t}\,B'$ with $\varsigma=\sqrt{128/30}$, and for $8T\sqrt Q$ in the terminal level. The base is $b=\lfloor(2\sqrt{Q/((n+\kappa)d)})^{1/2}\rceil$, which balances the first two terms of~\eqref{eq:bprime}. The projection is coded with $\sigma=\sqrt{(1+t)/2}\,B$ and the opening of the terminal level with $\sigma=\sqrt{Q/(nd)}$, so that by \cref{lem:coding} the sizes of the model are upper bounds for all proofs. We use $\theta=2$ and $t=128^{-1/2}$.

\paragraph{Search.}
The recursion is found by a beam search of width $160$ over at most eight levels; a state is a pair $(m,\beta)$, and the options of a level are $r$ from a geometric grid, decomposition or not, committed or clear partial evaluations, and the budget for the inner low parts. LaBRADOR is modelled along the same lines with the formulas of~\cite[Section~5]{BS23}, including the garbage polynomials and their decomposition, and its recursion is optimised in the same way.

\paragraph{First fold.}
The witness of~\eqref{eq:gh-system} has $\delta_wr+\delta_tnr+2\delta_0m+2n_1+2n$ ring elements, where $\delta_w$ and $\delta_t$ are the numbers of digits of $\vw$ and of the inner commitments. In the conventional and proof-based settings the shape $(m,r,b_0,\delta_0,b_1)$ is that of~\cite[Table~4]{NS24}, and the norm bound is computed from the variances of the digits, of $\vz$ and of the low parts as above; the inner commitments drop $\varrho_t\in\{8,9\}$ bits, which removes one digit, and the outer ones $\varrho_u\in\{11,12\}$ bits. In the guaranteed setting $m$ and $r$ are kept, the polynomial is split into $\delta_0=\lceil32/\log_2b_0\rceil$ digits for $b_0\in\{8,16\}$, the commitments have digits in base $2^6$, $\varrho_t$ is the largest number with $rnd\,4^{\varrho_t-1}\le B_s^2/10$, and the norm bound is computed as in \cref{sec:completeness}; the model selected $b_0=16$ with $\varrho_t\in\{7,8\}$ and $\varrho_u\in\{10,11\}$.

\section{Challenge Pairs for Proposition~\ref{prop:weak}}\label{app:pairs}

The following two challenges $c\in\Ch$ are such that the polynomial $c'$ obtained by exchanging the coefficients of $X^0$ and $X^1$ is also in $\Ch$; the differences are $c-c'=3(1-X)$ for the first and $4(1-X)$ for the second. Coefficients are listed from $X^0$ to $X^{63}$; the operator norms of $c$ and $c'$ are given after each list.

\input{gen/tab_pairs.tex}

\section{End-to-End Runs: Levels, Time and Memory}\label{app:e2e}

\Cref{tab:e2e-levels} lists, for $N=\nEeNlv$ and all runs of \cref{tab:e2e} at that size in which a proof was produced, the norm of the witness of every level relative to its bound and the numbers of attempts. In the proof-based setting the first level carries the variation of the first-fold witness, and the terminal level repeats its challenges until $\vz$ and $\ve$ meet their bounds; in the guaranteed setting every witness stays near half of its bound and one attempt sufficed almost always.

\begin{table}[h]
\centering\small
\setlength{\tabcolsep}{4pt}
\begin{tabular}{@{}lrrrrrr@{}}
\toprule
 & \multicolumn{3}{c}{proof-based} & \multicolumn{3}{c}{guaranteed}\\
\cmidrule(lr){2-4}\cmidrule(l){5-7}
level & $\norm{\vw}/\beta$ & $\Pi$ & $c$ & $\norm{\vw}/B$ & $\Pi$ & $c$\\
\midrule
\input{gen/tab_e2e_levels.tex}
\bottomrule
\end{tabular}
\caption{Runs of \cref{tab:e2e} for $N=\nEeNlv$: range of the norm of the witness of each level relative to its bound, and mean and maximum of the numbers of attempts for the projection ($\Pi$) and for the challenges ($c$); the row ``first fold'' gives the witness handed to the first level and the challenge vectors of the first fold.}
\label{tab:e2e-levels}
\end{table}

\Cref{tab:e2e-time} gives running times and memory. They were measured on one machine (\nEeMachine), with the prover and the verifier in separate processes, and include neither the generation of the public matrices nor the evaluation of $f$ that checks $y$.

\begin{table}[h]
\centering\small
\begin{tabular}{@{}llrrrrr@{}}
\toprule
 & & \multicolumn{3}{c}{time (s)} & \multicolumn{2}{c}{peak memory (MB)}\\
\cmidrule(lr){3-5}\cmidrule(l){6-7}
$N$ & setting & commit & open & verify & prover & verifier\\
\midrule
\input{gen/tab_e2e_time.tex}
\bottomrule
\end{tabular}
\caption{Median running times of the reference implementation for uniform polynomials (\nEeTimeNote) and peak resident memory of the prover (commitment and opening) and of the verifier.}
\label{tab:e2e-time}
\end{table}
